\chapter{Background and Related Work}
\label{ch:background}
% First prose draft, 2026-10-24. Mathematical foundations are attributed inline.

The statistical model, numerical method, and security decision in LM-IDNet
address different parts of the monitoring problem. A count distribution
describes the observations, an estimator fits its parameters, a numerical
backend evaluates the required functions, and a calibrated rule turns a
score into a decision. This chapter develops those distinctions before
comparing the system with related IoT detectors.

\section{IoT network anomaly detection}
\label{sec:background-detection}

An anomaly is an observation that departs from an expected pattern.
Chandola et al. organise anomaly-detection methods by their assumptions,
available labels, and means of distinguishing normal from unusual
observations \cite{chandola2009}. Statistical methods express expected
behaviour through a distribution and assess how compatible a new
observation is with that distribution. Their interpretation depends on
the adequacy of both the model and the observation supplied to it.

For a network detector, an observation may be a packet, a flow, or an
aggregate over a time interval. These are different units of analysis.
Packet-level observations can preserve ordering and short changes;
aggregate counts deliberately discard such detail. The choice determines
which deviations the model can represent. Two packet sequences that
produce the same count vector are indistinguishable to a detector whose
input is that vector alone.

Normal-reference learning estimates expected behaviour from traffic
assumed to be benign. Its practical advantage is that attack examples are
not required to fit the model. Its limitation is the need to justify the
reference data and its continuing relevance. Training contamination is
therefore a central issue: an anomaly sensor may learn activity that should
remain detectable. Cretu et al. address this problem through training-data
sanitisation \cite{cretu2008}. LM-IDNet instead assumes a trusted initial
reference and applies admission restrictions when considering later
learning data; the mechanisms should not be conflated.

A deviation also needs an operational interpretation. Unfamiliar
benign behaviour can generate false positives, and a malicious action may
preserve the features monitored by the detector. These limitations are
particularly relevant to the operational gap described by Sommer and
Paxson \cite{sommer2010}. In this thesis an alert reports statistical
evidence for investigation; attack attribution and response require
additional evidence.

The evaluation must preserve that distinction. Precision measures how many
positive decisions correspond to labelled attacks, whereas recall measures
how many labelled positives are detected. Their denominators differ, and
both matter when attacks are sparse. Saito and Rehmsmeier explain how
precision--recall analysis exposes effects of class imbalance that can be
obscured by other summaries \cite{saito2015}. A further distinction concerns
time: attacks can occupy intervals rather than isolated observations.
Tatbul et al. develop range-based precision and recall for this setting
\cite{tatbul2018}. The present work reports its implemented window and
attack-interval measures explicitly rather than treating them as
interchangeable.

\section{Categorical traffic counts and overdispersion}
\label{sec:background-counts}

\subsection{From packets to a count vector}

For a device and a completed observation window, let
\[
  \boldsymbol{x}=(x_1,\ldots,x_K),\qquad
  x_k\in\{0,1,2,\ldots\},\qquad
  N=\sum_{k=1}^{K}x_k.
\]
Here $K$ is the number of categories, $x_k$ the observed count in category
$k$, and $N$ the total number of counted packets. LM-IDNet uses $K=4$,
ordered as TCP, ordinary UDP, SSDP, and ARP. SSDP is separated from the
remaining UDP traffic, so a counted packet belongs to one category. Packets
outside the selected categories do not contribute to $N$. Consequently,
``total count'' refers to the total represented by this feature vector,
rather than every frame present in the source capture.

The four-category representation and ten-minute precedent come from the
IoT modelling study of Al-Haj Baddar et al. \cite{alhajbaddar2025iot}.
They constitute an application-specific representation. Other IoT studies use richer traffic
attributes: Sivanathan et al. investigate device characteristics, while
Doshi et al. select features using the communication properties of consumer
IoT devices \cite{sivanathan2019,doshi2018}. Their work motivates examining
device behaviour. This thesis evaluates the information retained by the
selected four-count representation.

Aggregation reduces the amount of data presented to the model and gives
each feature a direct interpretation. It also removes packet order,
inter-arrival structure within the window, endpoints, payload content, and
other information not encoded in the categories. The same protocol mixture
may therefore arise from different underlying activities. Window duration
adds another choice: a longer interval pools more observations but can
merge benign and attack traffic and postpone a decision until the window
closes. Duration must be evaluated against the intended detection and
response objective.

\subsection{The multinomial reference model}

Let $\boldsymbol{p}=(p_1,\ldots,p_K)$ be category probabilities, with
$p_k\geq 0$ and $\sum_k p_k=1$. Conditional on the total $N$, a multinomial
model assigns count vector $\boldsymbol{x}$ probability
\[
  P(\boldsymbol{X}=\boldsymbol{x}\mid N,\boldsymbol{p})
  =\frac{N!}{\prod_{k=1}^{K}x_k!}
   \prod_{k=1}^{K}p_k^{x_k}.
\]
For category $k$, its conditional mean and variance are
\[
  \operatorname{E}[X_k\mid N]=Np_k,
  \qquad
  \operatorname{Var}(X_k\mid N)=Np_k(1-p_k).
\]
The DM literature uses this fixed-probability model as the reference for
extra-multinomial variation \cite{yu2014,alhajbaddar2025analysis}.

A shared $\boldsymbol{p}$ permits random variation in the counts, but does
not allow the underlying category probabilities themselves to change
between windows. Ordinary device activity can introduce more variation
than this conditional model permits. In that case, overdispersion means
variation beyond the multinomial expectation at a comparable total and
mean profile. The relevant reference is the multinomial conditional
variance; the variance-to-mean criterion instead belongs to an equidispersed
Poisson model.

The conditioning on $N$ also matters when examining data. Raw counts from
windows with very different totals can have different variances even if
their category probabilities are unchanged. Assessing extra-multinomial
variation therefore requires a statistical comparison that retains the
observed totals.

\section{The Dirichlet--multinomial model}
\label{sec:background-dm}

\subsection{Compounding a multinomial with a Dirichlet distribution}

The compound-multinomial construction has a long statistical history;
Mosimann's 1962 paper is an early reference \cite{mosimann1962}.
The DM model introduces a random category-probability vector
$\boldsymbol{\theta}$ drawn from a Dirichlet distribution. Let
$\boldsymbol{\alpha}=(\alpha_1,\ldots,\alpha_K)$, where every
$\alpha_k>0$, and define the total concentration
$\alpha_0=\sum_{k=1}^{K}\alpha_k$. On the probability simplex, the
Dirichlet density is
\[
  f(\boldsymbol{\theta}\mid\boldsymbol{\alpha})
  =\frac{\Gamma(\alpha_0)}{\prod_{k=1}^{K}\Gamma(\alpha_k)}
    \prod_{k=1}^{K}\theta_k^{\alpha_k-1},
\]
where $\Gamma$ denotes Euler's Gamma function. A count window is modelled
hierarchically as
\[
  \boldsymbol{\theta}\sim\operatorname{Dirichlet}(\boldsymbol{\alpha}),
  \qquad
  \boldsymbol{X}\mid N,\boldsymbol{\theta}
    \sim\operatorname{Multinomial}(N,\boldsymbol{\theta}).
\]
Minka describes this compounding and its parameter estimation; Yu and
Shaw discuss the resulting count distribution and its numerical
evaluation \cite{minka2000,yu2014}.

Integrating out $\boldsymbol{\theta}$ gives the probability mass function
\begin{equation}
\begin{split}
  P(\boldsymbol{X}=\boldsymbol{x}\mid N,\boldsymbol{\alpha})
  ={}&\frac{N!}{\prod_{k=1}^{K}x_k!}
      \frac{\Gamma(\alpha_0)}{\Gamma(N+\alpha_0)}
      \prod_{k=1}^{K}
      \frac{\Gamma(x_k+\alpha_k)}{\Gamma(\alpha_k)}.
\end{split}
\label{eq:dm-pmf}
\end{equation}
The count support is defined by $\sum_k x_k=N$. The multinomial
coefficient in Eq.~\eqref{eq:dm-pmf} counts the orderings represented by a
count vector. It is needed when assigning probability to that vector,
while remaining constant with respect to $\boldsymbol{\alpha}$ during
maximisation on fixed observed data.

\subsection{Mean profile, concentration, and conditional variation}

The mean category profile is
$p_k=\operatorname{E}[\theta_k]=\alpha_k/\alpha_0$. Thus
$\boldsymbol{\alpha}=\alpha_0\boldsymbol{p}$ separates the mean mixture
from its concentration. The conditional first and second moments are
\begin{equation}
  \operatorname{E}[X_k\mid N]=Np_k,
  \qquad
  \operatorname{Var}(X_k\mid N)
    =Np_k(1-p_k)\frac{N+\alpha_0}{1+\alpha_0}.
\label{eq:dm-variance}
\end{equation}
These expressions follow by applying total expectation and total variance
to the Dirichlet--multinomial hierarchy
\cite{yu2014,alhajbaddar2025analysis}. For example, the variance combines
$\operatorname{E}[N\theta_k(1-\theta_k)]$ with
$\operatorname{Var}(N\theta_k)$; the Dirichlet variance
$p_k(1-p_k)/(1+\alpha_0)$ produces the factor in
Eq.~\eqref{eq:dm-variance}.

For fixed $\boldsymbol{p}$ and $N>1$, a smaller concentration permits
greater variation than the multinomial reference. For illustration,
$N=100$, $p_k=0.5$, and $\alpha_0=9$ give conditional variance $272.5$,
compared with $25$ for the multinomial. This algebraic example illustrates
the parameter effect. As $\alpha_0$ tends to
infinity at fixed $\boldsymbol{p}$, the inflation factor tends to one and
the DM approaches the multinomial model.

LM-IDNet reports the inverse concentration $\psi=1/\alpha_0$ as a
dispersion indicator. An alternative parameter is
$\rho=1/(1+\alpha_0)=\psi/(1+\psi)$, for which the inflation factor is
$1+(N-1)\rho$. Specifying the parameterisation is essential when comparing
reported dispersion values. Both $\psi$ and $\rho$ describe dispersion.

The categories are dependent because they share a total. For $j\ne k$,
\[
  \operatorname{Cov}(X_j,X_k\mid N)
  =-Np_jp_k\frac{N+\alpha_0}{1+\alpha_0}.
\]
This dependence is part of the count model. It should not be interpreted
as a learned causal relationship between protocols or as an arbitrary
model of positive correlation between their raw traffic volumes.

\subsection{What the model leaves unspecified}

The hierarchy models category composition conditional on the total. The
conditional mass still changes with $N$, because a different total changes
the count support. An explicit upper-tail count check therefore supplies
separate evidence about large totals. This distinction motivates the
complementary volume and ARP signals in the implemented detector.

When $N=0$, the only possible count vector is the zero vector, and
Eq.~\eqref{eq:dm-pmf} assigns it probability one. Availability monitoring
therefore requires a separate signal. Missing recording coverage remains a
separate data-quality state from an observed zero vector, and the input
contract distinguishes them before scoring.

Finally, the fitted DM is a distribution over windows. Using a product of
window probabilities as the fitting objective assumes independent rows
under the model. Adjacent traffic
windows can remain dependent in practice; seasonality, transitions, and
event order require a temporal extension.

\section{Parameter estimation and likelihood computation}
\label{sec:background-estimation}

\subsection{The fitting objective and the complete window score}

Suppose the fit set contains $M$ count vectors
$\boldsymbol{x}_1,\ldots,\boldsymbol{x}_M$, with totals
$N_i=\sum_k x_{ik}$. Taking logarithms of Eq.~\eqref{eq:dm-pmf} yields
\begin{equation}
\begin{split}
  \log P(\boldsymbol{x}_i\mid N_i,\boldsymbol{\alpha})
  ={}&\log\Gamma(N_i+1)-\sum_{k=1}^{K}\log\Gamma(x_{ik}+1)\\
     &+\log\Gamma(\alpha_0)-\log\Gamma(N_i+\alpha_0)\\
     &+\sum_{k=1}^{K}
       \bigl[\log\Gamma(x_{ik}+\alpha_k)-\log\Gamma(\alpha_k)\bigr].
\end{split}
\label{eq:dm-log-score}
\end{equation}
The first line depends only on the observed counts. Write it as
$C(\boldsymbol{x}_i)$ and denote the remaining parameter-dependent terms
by the kernel $\kappa(\boldsymbol{x}_i;\boldsymbol{\alpha})$. Then
\begin{equation}
  s_{\mathrm{DM}}(\boldsymbol{x}_i)
  =C(\boldsymbol{x}_i)+
    \kappa(\boldsymbol{x}_i;\boldsymbol{\alpha}).
\label{eq:kernel-and-score}
\end{equation}
The log-likelihood formulation and iterative estimation are discussed by
Minka and by Al-Haj Baddar et al.
\cite{minka2000,alhajbaddar2025analysis}.

For fixed training counts, maximising the sum of the complete
log-probabilities is equivalent to maximising
\[
  \mathcal{L}(\boldsymbol{\alpha})
  =\sum_{i=1}^{M}\kappa(\boldsymbol{x}_i;\boldsymbol{\alpha}),
\]
because $\sum_i C(\boldsymbol{x}_i)$ is constant with respect to
$\boldsymbol{\alpha}$. The same constant cancels when comparing the
objective at successive parameter iterates. This is the fitting contract
implemented by both likelihood backends.

Scoring a new count vector is a different operation. The coefficient
$C(\boldsymbol{x})$ varies between observations and must be restored to
obtain the complete log-probability in Eq.~\eqref{eq:kernel-and-score}.
Its omission could change the relative ordering of windows and the
calibration threshold, even when the fitted parameters remain identical.
Recorded fitting-kernel values and complete scoring values therefore have
different meanings.

Row-wise evaluation is equally important. In general,
\[
  \sum_{i=1}^{M}\kappa(\boldsymbol{x}_i;\boldsymbol{\alpha})
  \ne
  \kappa\!\left(\sum_{i=1}^{M}\boldsymbol{x}_i;
    \boldsymbol{\alpha}\right).
\]
Pooling counts before evaluating the nonlinear likelihood creates a
different observation. A matrix interface may reuse numerical setup
across rows, but must preserve the row-wise sum to retain this objective.

\subsection{Minka's fixed-point estimator}

The maximum-likelihood parameter estimate is obtained iteratively. Let
\[
  \digamma(z)\coloneqq\frac{\mathrm{d}}{\mathrm{d}z}\log\Gamma(z)
  =\frac{\Gamma'(z)}{\Gamma(z)}
\]
denote the digamma function.
The gradient component of the kernel sum is
\[
  \frac{\partial\mathcal{L}}{\partial\alpha_k}
  =\sum_{i=1}^{M}\bigl[
    \digamma(\alpha_0)-\digamma(N_i+\alpha_0)
    +\digamma(x_{ik}+\alpha_k)-\digamma(\alpha_k)\bigr].
\]
Minka gives the following fixed-point update for count data
\cite[Sec.~3, Eq.~(55)]{minka2000}:
\begin{equation}
  \alpha_k^{(r+1)}=\alpha_k^{(r)}
  \frac{\displaystyle\sum_{i=1}^{M}
    [\digamma(x_{ik}+\alpha_k^{(r)})-\digamma(\alpha_k^{(r)})]}
  {\displaystyle\sum_{i=1}^{M}
    [\digamma(N_i+\alpha_0^{(r)})-\digamma(\alpha_0^{(r)})]},
  \quad k=1,\ldots,K.
\label{eq:minka-update}
\end{equation}
Here $r$ indexes iterations and
$\alpha_0^{(r)}=\sum_k\alpha_k^{(r)}$. All components are updated using the
same preceding parameter vector.

Equation~\eqref{eq:minka-update} applies to count vectors. It is distinct
from estimating a Dirichlet density directly from observed proportions.
Replacing counts by normalised proportions and fitting that density would
discard the different totals and change the estimation problem. An
implementation also needs a positive initial vector, a stopping rule, and
an explicit outcome when a usable estimate is not reached. These are part
of the fitting procedure rather than properties supplied by the likelihood
backend.

In LM-IDNet, SciPy supplies the digamma evaluations for this update
\cite{virtanen2020}. The selected likelihood backend evaluates the
objective used to monitor convergence and later contributes to scoring.
Equation~\eqref{eq:minka-update} remains the estimator across backends. This
separation permits numerical implementations to be compared without
presenting them as different statistical estimators.

\subsection{From a score to a calibrated decision}

A complete log-probability supplies an ordering in which lower values
indicate observations assigned less mass by the normal model. An
elementary composition-only detector declares an anomaly when
$s_{\mathrm{DM}}(\boldsymbol{x})<\tau$, where $\tau$ is a cutoff estimated
from a separate benign calibration set. This is an instance of statistical
anomaly detection \cite{chandola2009}. LM-IDNet's retained rule combines
composition with count margins; its exact standardisation and joint
calibration are specified in Chapter~\ref{ch:implementation}.

Empirical quantiles allow a cutoff to be chosen without assuming Gaussian
scores, but their interpolation convention must be specified. Hyndman and
Fan compare the different sample-quantile definitions used in statistical
software \cite{hyndman1996}. The chosen quantile remains an empirical
operating point whose future false-positive rate depends on traffic
stability and dependence. Computing an attack posterior would additionally
require a distribution for attacks and their prior prevalence.

\section{Yu--Shaw and Languasco--Migliardi methods}
\label{sec:background-numerics}

\subsection{The numerical difficulty}

Logarithms avoid forming extremely small probability products directly,
while subtraction can still introduce floating-point error. In
Eq.~\eqref{eq:dm-log-score}, terms such as
$\log\Gamma(\alpha_k+x_k)-\log\Gamma(\alpha_k)$ are differences of
function values. Subtracting already rounded quantities can expose large
relative errors when the desired difference is small compared with the
operands. Goldberg discusses this cancellation problem and the value of
reformulating numerical expressions \cite{goldberg1991}.

For positive $a$ and a non-negative integer count $n$, define
\begin{equation}
  \Delta(a,n)=\log\Gamma(a+n)-\log\Gamma(a).
\label{eq:gamma-difference}
\end{equation}
Euler's gamma function satisfies the functional equation
$\Gamma(s+1)=s\Gamma(s)$ for $s>0$. Taking logarithms gives
$\log\Gamma(s+1)=\log s+\log\Gamma(s)$; repeated application yields
the exact identity
\[
  \Delta(a,n)=\sum_{j=0}^{n-1}\log(a+j)
  \quad\text{for }n\geq 1,
  \qquad \Delta(a,0)=0.
\]
This expresses the increment without subtracting two independently
computed log-Gamma values, but direct summation requires a number of terms
that grows with the count. It illustrates the accuracy--cost problem
addressed by specialised DM likelihood methods
\cite{yu2014,languasco2023}.

The kernel can be written using Eq.~\eqref{eq:gamma-difference} as
\[
  \kappa(\boldsymbol{x};\boldsymbol{\alpha})
    =\sum_{k=1}^{K}\Delta(\alpha_k,x_k)
      -\Delta(\alpha_0,N).
\]
Accurate evaluation of these increments is useful near the multinomial
limit, where the concentration can be large. Numerical checks therefore
cover the complete score, including its count-dependent coefficient.

\subsection{The Yu--Shaw approach and the LM method}

Yu and Shaw develop a method for accurate DM log-likelihood evaluation,
motivated by numerical instability near the multinomial limit and the
cost of alternative formulations for large counts \cite{yu2014}. Although
their application concerns biological count data, the likelihood
calculation is relevant to categorical traffic counts. The application
domain changes the observations, not the underlying special-function
problem.

Languasco and Migliardi show that the Yu--Shaw derivation is not
mathematically valid as written: it applies an asymptotic log-Gamma formula
without retaining either its large-argument condition or its asymptotic
error term. They also correct its truncation-error estimate and identify
error accumulation across the mesh computation, which Yu and Shaw did not
analyse \cite[Introduction, p.~1997; Sec.~4]{languasco2023}. The series
identity used by Yu and Shaw nevertheless remains correct; Languasco and
Migliardi supply a rigorous proof based on the Taylor expansion of
$\log(1+u)$.

Languasco and Migliardi subsequently introduced a rigorously proved method
for computing nearby log-Gamma differences with explicit error bounds and
a complete computational-cost analysis \cite{languasco2023}. Their
construction uses analytic reformulation, including Euler--Maclaurin
expansion and argument shifts, to evaluate the required differences. Its
error bounds account for accumulation across the likelihood terms and give
an a priori criterion for bounding the final numerical error in the DM log
probability of Eq.~\eqref{eq:dm-log-score}. In this thesis, \emph{LM} names
that numerical method; Minka's iteration remains the estimator, and
LM-IDNet supplies the anomaly decision.
Zero counts correspond to the exact zero increment in
Eq.~\eqref{eq:gamma-difference}, an important case for sparse protocol
windows.

Al-Haj Baddar et al. evaluate LM and Yu--Shaw within fixed-point DM
parameter estimation on count datasets from several domains
\cite{alhajbaddar2025analysis}. Their study supplies empirical evidence
about the numerical procedures under its chosen workloads. This thesis
measures SciPy timing and detector decisions under its own workloads.

The distinction determines the comparison in this thesis. The statistical
objective and count-score convention are kept fixed; numerical agreement,
decision agreement, and runtime are examined as separate outcomes. Matrix
fitting calls and isolated scoring calls have different amounts of work
and interface overhead. Their timings must be attached to the workload
measured, rather than combined into a single claim that a backend is
universally faster.

\section{Related systems and research gap}
\label{sec:background-related}

\subsection{Traffic profiles and IoT attack monitoring}

Sivanathan et al. use traffic characteristics to classify IoT devices
\cite{sivanathan2019}. A later study constructs one-class clustering
models and examines changes in device behaviour, including changes
associated with legitimate updates \cite{sivanathan2020}. Together, these
works support device-specific modelling while showing why a deviation
requires security interpretation. LM-IDNet instead uses a smaller
observation vector and a parametric count distribution.

Hamza et al. combine Manufacturer Usage Description (MUD) profiles with
software-defined networking (SDN) telemetry to monitor volumetric attacks
\cite{hamza2019}. Their approach observes activity associated with
permitted communication and includes attack traces released with the
study. This is relevant both to monitoring traffic volume and to the
provenance of the UNSW data used here. LM-IDNet consumes selected packet
counts within a separate statistical architecture.

Doshi et al. investigate machine-learning detection of consumer-IoT DDoS
traffic using features informed by IoT communication behaviour
\cite{doshi2018}. Their study provides another precedent for selecting
features with an intended network-middlebox setting in mind. It uses
different features and classifier experiments; the present work evaluates
its normal-only learning rule separately.

\subsection{Neural normal-profile detectors}

N-BaIoT uses deep autoencoders to identify departures from benign device
traffic and evaluates attacks involving Mirai and BASHLITE
\cite{meidan2018}. Its traffic statistics and learned representation differ
from LM-IDNet's four counts. It is a relevant comparison of modelling
choices; direct benchmark results here use the shared local data.

Kitsune combines incremental traffic-feature extraction with an ensemble
of small autoencoders, called KitNET, for online network intrusion
detection \cite{mirsky2018}. Its evaluation includes a Raspberry Pi
setting. This demonstrates that online and resource-conscious neural
detection already has an established precedent. LM-IDNet resource results
come from its own constrained simulation.

D\"IoT builds device-type communication profiles and aggregates learning
across gateways through a federated approach \cite{nguyen2019}. It models
communication sequences, whereas the DM detector treats completed count
windows as its statistical observations. The comparison concerns what
is represented, how normal behaviour is learned, and where computation is
placed.

Table~\ref{tab:related-systems} collects these architectural distinctions.
Detection accuracy remains within each study's traffic, attacks, observation
units, and evaluation protocol. The experimental baselines in this thesis
are the locally implemented variants described in Chapter~\ref{ch:experiments}.

\begin{table}[tbp]
\centering
\small
\begin{tabular}{@{}>{\raggedright\arraybackslash}p{0.19\textwidth}
  >{\raggedright\arraybackslash}p{0.34\textwidth}
  >{\raggedright\arraybackslash}p{0.40\textwidth}@{}}
\toprule
\textbf{System or study} & \textbf{Representation and learning} &
\textbf{Relation to this thesis} \\
\midrule
Sivanathan et al. \cite{sivanathan2020} &
Flow attributes and one-class clustering. &
Device profiles and behavioural change; a different feature and model choice. \\
Hamza et al. \cite{hamza2019} &
MUD activity and SDN telemetry. &
Volumetric monitoring and the associated UNSW attack-data provenance. \\
N-BaIoT \cite{meidan2018} &
Traffic statistics and deep autoencoders. &
Normal-device learning with a richer representation and different attack data. \\
Kitsune \cite{mirsky2018} &
Incremental features and an autoencoder ensemble. &
Online gateway-oriented detection with its own resource measurements. \\
D\"IoT \cite{nguyen2019} &
Communication sequences and federated device-type models. &
Distributed profile learning compared with local count-window modelling. \\
LM-IDNet (this work) &
Four count categories, a DM model, and complementary count margins. &
Normal-reference calibration, LM numerical integration, and guarded updates
evaluated in a common pipeline. \\
\bottomrule
\end{tabular}
\caption{Qualitative comparison of related systems with LM-IDNet. The entries
describe architectural choices; different datasets and evaluation units
prevent interpreting the table as a detection-performance ranking.}
\label{tab:related-systems}
\end{table}

\subsection{Drift, poisoning, and the interpretation of adaptation}

Gama et al. discuss concept drift primarily in an online supervised
setting, where the relationship between observations and a target changes
\cite{gama2014}. The present detector learns from unlabelled runtime
observations. Changes in those inputs are therefore described
as behavioural or distributional change unless the stronger concept-drift
interpretation is justified. Adaptation quality is measured through its
effect on later outcomes.

Kloft and Laskov analyse how admitted observations can poison an online
centroid anomaly detector and how constraints affect the attacker's
influence \cite{kloft2012}. Their formal results apply to the analysed
centroid model. They motivate examining admission and reference constraints
in LM-IDNet. Cretu et al.'s sanitisation work similarly motivates explicit
measurement of contamination risk \cite{cretu2008}. Security bounds would
require analysis of the specific DM update rule.

There is an inherent tension in retaining an original normal reference.
It can preserve evidence against activity that a moving model might
otherwise accept. It can also preserve an alert caused by a legitimate
change. The adaptive evaluation must therefore ask which observations were
admitted, what changed, and whether the change improved a relevant outcome.
If no candidate is promoted, the experiment documents the behaviour of the
guards and the continued use of the starting model.

\subsection{The question addressed by LM-IDNet}

The closest statistical precursor is the DM IoT study of Al-Haj Baddar et
al. \cite{alhajbaddar2025iot}. The present work retains its count-modelling
foundation and examines the additional decisions needed for a detector:
separate benign calibration, a later labelled evaluation stream,
complementary count checks, candidate admission and activation, and
traceable event and state records. These are the particular integration
and evaluation questions addressed here.

The experimental design is also part of the question. TESSERACT shows how
temporal and distributional bias can distort a security-classification
evaluation \cite{pendlebury2019}, while Arp et al. discuss pitfalls such as
data snooping, unsuitable measures, and lab-only conclusions
\cite{arp2022}. Applied to this study, these concerns require disclosure
of development choices, the provenance of labels, chronological ordering,
and the boundary between the local simulation and physical deployment.
The following chapters define the implemented detector and the evidence
used to answer those questions.
